\documentclass[10pt,a4paper]{amsart}
\usepackage{amsmath,amssymb,amsthm,graphicx}
\usepackage{verbatim}
\newenvironment{reference}{\comment}{\endcomment}
%\newenvironment{reference}{}{}
\newenvironment{slogan}{\comment}{\endcomment}
\newenvironment{history}{\comment}{\endcomment}
% For commutative diagrams we use Xy-pic
\usepackage[all]{xy}

% We use 2cell for 2-commutative diagrams.
\xyoption{2cell}
% Theorem environments.
%
\theoremstyle{plain}
\newtheorem{theorem}[subsection]{Theorem}
\newtheorem{proposition}[subsection]{Proposition}
\newtheorem{lemma}[subsection]{Lemma}

\theoremstyle{definition}
\newtheorem{definition}[subsection]{Definition}
\newtheorem{example}[subsection]{Example}
\newtheorem{exercise}[subsection]{Exercise}
\newtheorem{situation}[subsection]{Situation}

\theoremstyle{remark}
\newtheorem{remark}[subsection]{Remark}
\newtheorem{remarks}[subsection]{Remarks}

\numberwithin{equation}{subsection}

% Macros
%
\def\lim{\mathop{\mathrm{lim}}\nolimits}
\def\colim{\mathop{\mathrm{colim}}\nolimits}
\def\Spec{\mathop{\mathrm{Spec}}}
\def\Hom{\mathop{\mathrm{Hom}}\nolimits}
\def\Ext{\mathop{\mathrm{Ext}}\nolimits}
\def\SheafHom{\mathop{\mathcal{H}\!\mathit{om}}\nolimits}
\def\SheafExt{\mathop{\mathcal{E}\!\mathit{xt}}\nolimits}
\def\Sch{\mathit{Sch}}
\def\Mor{\mathop{\mathrm{Mor}}\nolimits}
\def\Ob{\mathop{\mathrm{Ob}}\nolimits}
\def\Sh{\mathop{\mathit{Sh}}\nolimits}
\def\NL{\mathop{N\!L}\nolimits}
\def\CH{\mathop{\mathrm{CH}}\nolimits}
\def\proetale{{pro\text{-}\acute{e}tale}}
\def\etale{{\acute{e}tale}}
\def\QCoh{\mathit{QCoh}}
\def\Ker{\mathop{\mathrm{Ker}}}
\def\Im{\mathop{\mathrm{Im}}}
\def\Coker{\mathop{\mathrm{Coker}}}
\def\Coim{\mathop{\mathrm{Coim}}}

% Boxtimes
%
\DeclareMathSymbol{\boxtimes}{\mathbin}{AMSa}{"02}

%
% Macros for moduli stacks/spaces
%
\def\QCohstack{\mathcal{QC}\!\mathit{oh}}
\def\Cohstack{\mathcal{C}\!\mathit{oh}}
\def\Spacesstack{\mathcal{S}\!\mathit{paces}}
\def\Quotfunctor{\mathrm{Quot}}
\def\Hilbfunctor{\mathrm{Hilb}}
\def\Curvesstack{\mathcal{C}\!\mathit{urves}}
\def\Polarizedstack{\mathcal{P}\!\mathit{olarized}}
\def\Complexesstack{\mathcal{C}\!\mathit{omplexes}}
% \Pic is the operator that assigns to X its picard group, usage \Pic(X)
% \Picardstack_{X/B} denotes the Picard stack of X over B
% \Picardfunctor_{X/B} denotes the Picard functor of X over B
\def\Pic{\mathop{\mathrm{Pic}}\nolimits}
\def\Picardstack{\mathcal{P}\!\mathit{ic}}
\def\Picardfunctor{\mathrm{Pic}}
\def\Deformationcategory{\mathcal{D}\!\mathit{ef}}
\usepackage{cleveref}
\crefname{originallemma}{original lemma}{original lemmas}
\Crefname{originallemma}{Original lemma}{Original lemmas}
\makeatletter
\newlabel{coherent-lemma-quasi-coherent-affine-cohomology-zero}{{01XB}{source}{}{sourcepointer-1}{}}
\newlabel{coherent-lemma-quasi-coherent-affine-cohomology-zero@cref}{{[originallemma][0][]01XB}{[1][0][]source}{}{}{}}
\makeatother
\crefname{originallemma}{original lemma}{original lemmas}
\Crefname{originallemma}{Original lemma}{Original lemmas}
\makeatletter
\newlabel{cohomology-lemma-cech-complex-complex-computes}{{0FLH}{source}{}{sourcepointer-2}{}}
\newlabel{cohomology-lemma-cech-complex-complex-computes@cref}{{[originallemma][0][]0FLH}{[1][0][]source}{}{}{}}
\makeatother
\crefname{originalsection}{original section}{original sections}
\Crefname{originalsection}{Original section}{Original sections}
\makeatletter
\newlabel{cohomology-section-cech-cohomology-of-complexes}{{01FP}{source}{}{sourcepointer-3}{}}
\newlabel{cohomology-section-cech-cohomology-of-complexes@cref}{{[originalsection][0][]01FP}{[1][0][]source}{}{}{}}
\makeatother
\crefname{originalsection}{original section}{original sections}
\Crefname{originalsection}{Original section}{Original sections}
\makeatletter
\newlabel{cohomology-section-cup-product}{{0FKU}{source}{}{sourcepointer-4}{}}
\newlabel{cohomology-section-cup-product@cref}{{[originalsection][0][]0FKU}{[1][0][]source}{}{}{}}
\makeatother
\crefname{originallemma}{original lemma}{original lemmas}
\Crefname{originallemma}{Original lemma}{Original lemmas}
\makeatletter
\newlabel{lemma-de-rham-affine}{{0FLW}{source}{}{sourcepointer-5}{}}
\newlabel{lemma-de-rham-affine@cref}{{[originallemma][0][]0FLW}{[1][0][]source}{}{}{}}
\makeatother
\crefname{originallemma}{original lemma}{original lemmas}
\Crefname{originallemma}{Original lemma}{Original lemmas}
\makeatletter
\newlabel{lemma-relative-global-generation-on-fibres}{{0FMP}{source}{}{sourcepointer-6}{}}
\newlabel{lemma-relative-global-generation-on-fibres@cref}{{[originallemma][0][]0FMP}{[1][0][]source}{}{}{}}
\makeatother
\crefname{originallemma}{original lemma}{original lemmas}
\Crefname{originallemma}{Original lemma}{Original lemmas}
\makeatletter
\newlabel{lemma-spectral-sequence-smooth}{{0FMM}{source}{}{sourcepointer-7}{}}
\newlabel{lemma-spectral-sequence-smooth@cref}{{[originallemma][0][]0FMM}{[1][0][]source}{}{}{}}
\makeatother
\crefname{originallemma}{original lemma}{original lemmas}
\Crefname{originallemma}{Original lemma}{Original lemmas}
\makeatletter
\newlabel{perfect-lemma-K-flat}{{0FMQ}{source}{}{sourcepointer-8}{}}
\newlabel{perfect-lemma-K-flat@cref}{{[originallemma][0][]0FMQ}{[1][0][]source}{}{}{}}
\makeatother
\title{Classes whose restrictions form a basis on every fibre make the de Rham cohomology of the total space free over that of the base}
\author{The Stacks Project Authors}
\date{stacks-project:derham.tex:local-snapshot-20261012; GNU FDL 1.2 or later}
\begin{document}
\maketitle
\noindent\textit{Editorial note.} All delivered mathematical text is the author's own, copied from the pinned Stacks Project chapter file derham.tex; no statement or proof step was rewritten, reordered or removed.\par\medskip
\noindent\textit{Editorial note.} Preamble provenance: the pinned chapter file opens with \textbackslash{}input\{preamble\} and carries no preamble text of its own, so the project's single Stacks-wide preamble staged at sources/stacks-derham/preamble.tex is excerpted verbatim as three spans, lines 8-12 (verbatim package and the project's own reference/slogan/history display environments), lines 14-18 (Xy-pic setup) and lines 154-217 (theorem declarations and mathematical macros). The wrapper class and the amsmath/amssymb/amsthm/graphicx base packages are supplied by the assembler. No macro was invented or redefined.\par\medskip
\noindent\textit{Editorial note.} Same-chapter pointers are located at their \textbackslash{}label lines inside the pinned chapter file; cross-chapter pointers are located in the project's own tag list staged in the snapshot (sources/stacks-chow/tags.txt), at the verbatim lines carrying their canonical tags.\par\medskip
\noindent\textit{Allowed standard prerequisites (editorial).}\begin{itemize}
\item The affine case of the de Rham complex computing its own derived global sections. Same chapter derham.tex, label lemma-de-rham-affine (invoked at proof line 2046). It identifies \$\textbackslash{}Omega\textasciicircum{}\textbackslash{}bullet\_\{A/R\}\$ with the derived global sections of the de Rham complex of \$Y\$ over \$S\$, which is the source of the map under study.
\item The Cech resolution of a complex by a finite affine open covering, and the facts used about the resulting differential graded algebra. Cohomology, Section \textbackslash{}ref\{cohomology-section-cech-cohomology-of-complexes\} (invoked at proof line 2054), labels cohomology-lemma-cech-complex-complex-computes (invoked at proof line 2063) and Section \textbackslash{}ref\{cohomology-section-cup-product\} (invoked at proof line 2072), Cohomology of Schemes, label coherent-lemma-quasi-coherent-affine-cohomology-zero (invoked at proof line 2061), and Same chapter derham.tex, label lemma-spectral-sequence-smooth (invoked at proof line 2092). They supply the complex \$K\textasciicircum{}\textbackslash{}bullet = \textbackslash{}operatorname\{Tot\}(\textbackslash{}check\{\textbackslash{}mathcal\{C\}\}\textasciicircum{}\textbackslash{}bullet(\textbackslash{}mathcal\{U\}, \textbackslash{}Omega\textasciicircum{}\textbackslash{}bullet\_\{X/S\}))\$ representing \$R\textbackslash{}Gamma(X, \textbackslash{}Omega\textasciicircum{}\textbackslash{}bullet\_\{X/S\})\$, its multiplication inducing the cup product, and the filtered structure whose associated graded pieces are the objects compared below; the author states these standard facts without proof and refers to the discussion before the spectral sequence lemma.
\item The relative form of the basis hypothesis on one fibre. Same chapter derham.tex, label lemma-relative-global-generation-on-fibres (invoked at proof line 2117). It shows that the induced map on the first associated graded piece, \$\textbackslash{}operatorname\{gr\}\textasciicircum{}0\textbackslash{}tilde x\$, is an isomorphism in \$D(A)\$.
\item K-flatness of the Cech complex, making the tensor product on the right hand side derived. Derived Categories of Schemes, label perfect-lemma-K-flat (invoked at proof line 2137). It lets the comparison of the higher graded pieces be obtained from the case \$k = 0\$ by derived tensoring with the free module \$\textbackslash{}Omega\textasciicircum{}k\_\{A/R\}[-k]\$.
\end{itemize}\medskip
\ifdefined\hypertarget\hypertarget{sourcepointer-1}{}\fi
\noindent\textit{Source reference.} Original lemma 01XB: sources/stacks-chow/tags.txt, source lines 2512--2512. Navigation pointer for \textbackslash{}ref\{coherent-lemma-quasi-coherent-affine-cohomology-zero\} used inside the counted span (line 2062); the argument that consumes it is retained in the span.\par\medskip
\ifdefined\hypertarget\hypertarget{sourcepointer-2}{}\fi
\noindent\textit{Source reference.} Original lemma 0FLH: sources/stacks-chow/tags.txt, source lines 19497--19497. Navigation pointer for \textbackslash{}ref\{cohomology-lemma-cech-complex-complex-computes\} used inside the counted span (line 2063); the argument that consumes it is retained in the span.\par\medskip
\ifdefined\hypertarget\hypertarget{sourcepointer-3}{}\fi
\noindent\textit{Source reference.} Original section 01FP: sources/stacks-chow/tags.txt, source lines 1878--1878. Navigation pointer for \textbackslash{}ref\{cohomology-section-cech-cohomology-of-complexes\} used inside the counted span (line 2054); the argument that consumes it is retained in the span.\par\medskip
\ifdefined\hypertarget\hypertarget{sourcepointer-4}{}\fi
\noindent\textit{Source reference.} Original section 0FKU: sources/stacks-chow/tags.txt, source lines 19474--19474. Navigation pointer for \textbackslash{}ref\{cohomology-section-cup-product\} used inside the counted span (line 2072); the argument that consumes it is retained in the span.\par\medskip
\ifdefined\hypertarget\hypertarget{sourcepointer-5}{}\fi
\noindent\textit{Source reference.} Original lemma 0FLW: sources/stacks-derham/pinned-original-source.tex, source lines 376--376. Navigation pointer for \textbackslash{}ref\{lemma-de-rham-affine\} used inside the counted span (line 2046); the argument that consumes it is retained in the span.\par\medskip
\ifdefined\hypertarget\hypertarget{sourcepointer-6}{}\fi
\noindent\textit{Source reference.} Original lemma 0FMP: sources/stacks-derham/pinned-original-source.tex, source lines 1987--1987. Navigation pointer for \textbackslash{}ref\{lemma-relative-global-generation-on-fibres\} used inside the counted span (line 2117); the argument that consumes it is retained in the span.\par\medskip
\ifdefined\hypertarget\hypertarget{sourcepointer-7}{}\fi
\noindent\textit{Source reference.} Original lemma 0FMM: sources/stacks-derham/pinned-original-source.tex, source lines 1911--1911. Navigation pointer for \textbackslash{}ref\{lemma-spectral-sequence-smooth\} used inside the counted span (line 2092); the argument that consumes it is retained in the span.\par\medskip
\ifdefined\hypertarget\hypertarget{sourcepointer-8}{}\fi
\noindent\textit{Source reference.} Original lemma 0FMQ: sources/stacks-chow/tags.txt, source lines 19540--19540. Navigation pointer for \textbackslash{}ref\{perfect-lemma-K-flat\} used inside the counted span (line 2137); the argument that consumes it is retained in the span.\par\medskip
\begin{lemma}
\label{lemma-global-generation-on-fibres}
Let $f : X \to Y$ be a smooth proper morphism of schemes over a base $S$.
Assume
\begin{enumerate}
\item $Y$ and $S$ are affine, and
\item there exist integers $N$ and $n_1, \ldots, n_N \geq 0$ and
$\xi_i \in H^{n_i}_{dR}(X/S)$, $1 \leq i \leq N$ such that
for all points $y \in Y$ the images of $\xi_1, \ldots, \xi_N$
in $H^*_{dR}(X_y/y)$ form a basis over $\kappa(y)$.
\end{enumerate}
Then the map
$$
\bigoplus\nolimits_{i = 1}^N H^*_{dR}(Y/S) \longrightarrow
H^*_{dR}(X/S), \quad
(a_1, \ldots, a_N) \longmapsto  \sum \xi_i \cup f^*a_i
$$
is an isomorphism.
\end{lemma}

\begin{proof}
Say $Y = \Spec(A)$ and $S = \Spec(R)$.
In this case $\Omega^\bullet_{A/R}$ computes
$R\Gamma(Y, \Omega^\bullet_{Y/S})$ by Lemma \ref{lemma-de-rham-affine}.
Choose a finite affine open covering $\mathcal{U} : X = \bigcup_{i \in I} U_i$.
Consider the complex
$$
K^\bullet =
\text{Tot}(\check{\mathcal{C}}^\bullet(\mathcal{U}, \Omega_{X/S}^\bullet))
$$
as in
Cohomology, Section \ref{cohomology-section-cech-cohomology-of-complexes}.
Let us collect some facts about this complex most of which
can be found in the reference just given:
\begin{enumerate}
\item $K^\bullet$ is a complex of $R$-modules whose terms are
$A$-modules,
\item $K^\bullet$ represents $R\Gamma(X, \Omega^\bullet_{X/S})$ in $D(R)$
(Cohomology of Schemes, Lemma
\ref{coherent-lemma-quasi-coherent-affine-cohomology-zero} and
Cohomology, Lemma \ref{cohomology-lemma-cech-complex-complex-computes}),
\item there is a natural map $\Omega^\bullet_{A/R} \to K^\bullet$
of complexes of $R$-modules which is $A$-linear on terms and
induces the pullback map $H^*_{dR}(Y/S) \to H^*_{dR}(X/S)$
on cohomology,
\item $K^\bullet$ has a multiplication denoted $\wedge$
which turns it into a differential graded $R$-algebra,
\item the multiplication on $K^\bullet$
induces the cup product on $H^*_{dR}(X/S)$
(Cohomology, Section \ref{cohomology-section-cup-product}),
\item the filtration $F$ on $\Omega^*_{X/S}$ induces a filtration
$$
K^\bullet =
F^0K^\bullet \supset F^1K^\bullet \supset F^2K^\bullet \supset \ldots
$$
by subcomplexes on $K^\bullet$ such that
\begin{enumerate}
\item $F^kK^n \subset K^n$ is an $A$-submmodule,
\item $F^kK^\bullet \wedge F^lK^\bullet \subset F^{k + l}K^\bullet$,
\item $\text{gr}^kK^\bullet$ is a complex of $A$-modules,
\item $\text{gr}^0K^\bullet =
\text{Tot}(\check{\mathcal{C}}^\bullet(\mathcal{U}, \Omega_{X/Y}^\bullet))$
and represents $R\Gamma(X, \Omega^\bullet_{X/Y})$ in $D(A)$,
\item multiplication induces an isomorphism
$\Omega^k_{A/R}[-k] \otimes_A \text{gr}^0K^\bullet \to \text{gr}^kK^\bullet$
\end{enumerate}
\end{enumerate}
We omit the detailed proofs of these statements; please see discussion
leading up to the construction of the spectral sequence in
Lemma \ref{lemma-spectral-sequence-smooth}.

\medskip\noindent
For every $i = 1, \ldots, N$ we choose a cocycle $x_i \in K^{n_i}$
representing $\xi_i$. Next, we look at the map of complexes
$$
\tilde x :
M^\bullet = \bigoplus\nolimits_{i = 1, \ldots, N}
\Omega^\bullet_{A/R}[-n_i]
\longrightarrow
K^\bullet
$$
which sends $\omega$ in the $i$th summand to $x_i \wedge \omega$.
All that remains is to show that this map is a quasi-isomorphism.
We endow $M^\bullet$ with the structure of a filtered complex
by the rule
$$
F^kM^\bullet =
\bigoplus\nolimits_{i = 1, \ldots, N}
(\sigma_{\geq k}\Omega^\bullet_{A/R})[-n_i]
$$
With this choice the map $\tilde x$ is a morphism of filtered complexes.
Observe that $\text{gr}^0M^\bullet = \bigoplus A[-n_i]$
and multiplication induces an isomorphism
$\Omega^k_{A/R}[-k] \otimes_A \text{gr}^0M^\bullet \to \text{gr}^kM^\bullet$.
By construction and Lemma \ref{lemma-relative-global-generation-on-fibres}
we see that
$$
\text{gr}^0\tilde x :
\text{gr}^0M^\bullet \longrightarrow
\text{gr}^0K^\bullet
$$
is an isomorphism in $D(A)$. It follows that for all $k \geq 0$
we obtain isomorphisms
$$
\text{gr}^k \tilde x :
\text{gr}^kM^\bullet = \Omega^k_{A/R}[-k] \otimes_A \text{gr}^0M^\bullet
\longrightarrow
\Omega^k_{A/R}[-k] \otimes_A \text{gr}^0K^\bullet =
\text{gr}^kK^\bullet
$$
in $D(A)$. Namely, the complex
$\text{gr}^0K^\bullet =
\text{Tot}(\check{\mathcal{C}}^\bullet(\mathcal{U}, \Omega_{X/Y}^\bullet))$
is K-flat as a complex of $A$-modules by Derived Categories of Schemes,
Lemma \ref{perfect-lemma-K-flat}.
Hence the tensor product on the right hand side is the
derived tensor product as is true by inspection on the left hand side.
Finally, taking the derived tensor product
$\Omega^k_{A/R}[-k] \otimes_A^\mathbf{L} -$ is a functor on $D(A)$
and therefore sends isomorphisms to isomorphisms.
Arguing by induction on $k$ we deduce that
$$
\tilde x : M^\bullet/F^kM^\bullet \to K^\bullet/F^kK^\bullet
$$
is an isomorphism in $D(R)$ since we have the short exact sequences
$$
0 \to F^kM^\bullet/F^{k + 1}M^\bullet \to
M^\bullet/F^{k + 1}M^\bullet \to
\text{gr}^kM^\bullet \to 0
$$
and similarly for $K^\bullet$. This proves that $\tilde x$ is a
quasi-isomorphism as the filtrations are finite in any given degree.
\end{proof}

\end{document}
